\documentclass[11pt,a4paper]{article}

\usepackage[margin=1in]{geometry}
\usepackage[T1]{fontenc}
\usepackage[utf8]{inputenc}
\usepackage{lmodern}
\usepackage{setspace}
\usepackage{graphicx}
\usepackage{booktabs}
\usepackage{longtable}
\usepackage{array}
\usepackage{hyperref}
\usepackage{enumitem}
\usepackage{pdfpages}
\usepackage{float}

\hypersetup{
    colorlinks=true,
    linkcolor=black,
    citecolor=black,
    urlcolor=blue
}

\setstretch{1.0}
\setlength{\parskip}{0.5em}
\setlength{\parindent}{0pt}

\title{ERRGen: Evidence-Grounded Equity Research Report Generator}
\author{Your Name \\ Student ID \\ Course / Project Code}
\date{\today}

\begin{document}

% ------------------------------------------------------------------
% Cover sheet
% ------------------------------------------------------------------
% Recommended workflow:
% 1. Convert the Word cover sheet to PDF.
% 2. Upload it to Overleaf as coversheet.pdf.
% 3. Uncomment the next line.
% 4. Comment out \maketitle, \thispagestyle{empty}, and the following \newpage.
%
% \includepdf[pages=-]{coversheet.pdf}

\maketitle
\thispagestyle{empty}
\newpage

\begin{abstract}
This project studies the problem of generating equity research reports with
strong factual grounding, numerical correctness, and traceable evidence.
Traditional large language model based report generation can produce fluent
text, but it often suffers from unsupported claims, numerical mistakes, weak
citations, and temporal leakage. To address these issues, we developed
ERRGen, an evidence-grounded agentic pipeline for professional equity research
report generation.

ERRGen retrieves data from real financial and news sources, converts them into
structured evidence chunks, performs deterministic financial calculations, and
generates section-level analysis with paragraph-level citations. Each generated
paragraph is then passed through a checker-reviser verification loop before it
can be accepted into the final report. In addition, the final investment
recommendation is gated so that it is skipped when unresolved major issues
remain in the upstream sections.

Experimental results show that the full ERRGen pipeline outperforms an
unverified baseline in pairwise evaluation, achieving 12 wins, 4 losses, and
4 ties over 20 samples, corresponding to an adjusted win rate of 0.70.
However, the generated reports still underperform human-written ground-truth
reports in overall analytical quality. These results suggest that verification
and evidence grounding improve reliability, while deeper reasoning and richer
financial synthesis remain important future directions.
\end{abstract}

\newpage
\tableofcontents
\newpage

\section{Introduction}

\subsection{Background and Motivation}
Large language models are increasingly capable of generating long-form business
and financial documents. However, equity research is a high-stakes setting in
which fluent writing alone is insufficient. A useful research report must be
factually grounded, numerically correct, temporally valid, and transparent
about the evidence supporting each conclusion.

In practice, a standard prompt-only generation workflow can produce plausible
but unreliable outputs. Common failure modes include hallucinated company
developments, unsupported investment conclusions, wrong calculations, and the
use of information beyond the requested as-of date. These weaknesses motivate
the design of a system that treats evidence tracking and verification as
first-class components rather than optional post-processing steps.

\subsection{Problem Statement}
The goal of this project is to build a system that generates an equity
research report for a target company given a ticker, an as-of date, and a set
of focus areas. The system should produce a report that is useful to an
investor while maintaining strong factual reliability and auditability.

More specifically, the report should satisfy the following design goals:
\begin{itemize}[leftmargin=1.5em]
    \item analytical claims should be grounded in retrieved evidence;
    \item numerical statements should come from deterministic calculations
    rather than free-form language model arithmetic;
    \item citations should be traceable at the paragraph level;
    \item content should not use evidence beyond the requested as-of date;
    \item the final recommendation should not be produced if major upstream
    issues remain unresolved.
\end{itemize}

\subsection{Project Overview}
To address the above problem, we developed ERRGen, an evidence-grounded
equity research report generator. The system takes a natural-language user
request, retrieves multi-source financial and news data, extracts structured
facts, computes deterministic financial metrics, generates section-level
analysis, verifies each paragraph through an iterative checker-reviser loop,
and finally assembles a report with evidence and calculation appendices.

Unlike a one-shot text generation system, ERRGen is designed around explicit
control points for factuality. It combines retrieval, structured evidence
representation, deterministic financial reasoning, and iterative verification
into a single workflow.

\subsection{Main Contributions}
The main contributions of this project are as follows:
\begin{itemize}[leftmargin=1.5em]
    \item We propose an evidence-grounded agentic pipeline for long-form equity
    research report generation.
    \item We introduce paragraph-level citation tracking together with
    deterministic financial calculations for auditable numerical reasoning.
    \item We design a checker-reviser verification loop that identifies and
    corrects unsupported claims, citation mismatches, numerical errors, and
    temporal violations before the report is finalized.
    \item We implement a gated investment recommendation mechanism so that the
    final recommendation is withheld when upstream analysis remains unreliable.
    \item We build an evaluation workflow that compares the full system against
    both an ablation baseline and ground-truth reports.
\end{itemize}

\subsection{My Individual Contributions}
For a team project, this subsection should clearly state the parts completed by
you. Replace the placeholder items below with your actual work and estimated
share.

\begin{itemize}[leftmargin=1.5em]
    \item Data retrieval and API integration: XX\%
    \item Evidence extraction and deterministic calculator: XX\%
    \item Analysis agent implementation: XX\%
    \item Verification and revision loop: XX\%
    \item Evaluation scripts and experiments: XX\%
    \item Testing, debugging, and report writing: XX\%
\end{itemize}

You can also rewrite this subsection in first person, for example:
\textit{I was mainly responsible for implementing the verification loop,
building the evaluation pipeline, and running the experiments.}

\section{Overview of Related Works}

\subsection{LLM-Based Long-Form Report Generation}
This subsection should review prior work on large language model based
long-form text generation, especially systems that aim to produce structured
reports or professional documents. Focus on the strengths of these systems in
fluency and organization, while also noting that long-form generation often
struggles with factual stability over many paragraphs.

\subsection{Retrieval-Augmented and Evidence-Grounded Generation}
This subsection should discuss retrieval-augmented generation and evidence
grounding methods. The key connection to our project is that generated content
should be anchored to retrieved external evidence rather than relying entirely
on the model's parametric memory.

\subsection{Self-Verification and Iterative Revision}
This subsection should summarize work on self-correction, fact-checking, and
iterative revision in LLM systems. This is closely related to our
checker-reviser loop, where a draft paragraph is validated and revised before
being accepted.

\subsection{Evaluation with LLM-as-a-Judge}
This subsection should review how LLMs are used as evaluators for generated
outputs. You can explain that our project uses an LLM judge to evaluate both
accuracy-oriented and quality-oriented dimensions, which is useful for comparing
different system variants at scale.

\subsection{Positioning of Our Work}
Conclude this section by explaining how ERRGen differs from existing work. The
key point is that our system combines evidence-grounded generation,
deterministic numerical reasoning, paragraph-level citation tracking, and a
verification gate in the specific domain of equity research report generation.

\section{Methodology and Results}

\subsection{Task Definition and Design Goals}
The input to ERRGen is a user request specifying a company, ticker, as-of date,
and optional focus areas. The output is a structured equity research report
covering company overview, recent developments, financial analysis, business
and competitive analysis, risk analysis, and, when appropriate, an investment
recommendation and outlook.

The system is designed to optimize not only readability but also factual
grounding, numerical correctness, temporal validity, and auditability.

\subsection{System Overview}
Figure~\ref{fig:pipeline} illustrates the overall ERRGen workflow.

\begin{figure}[H]
\centering
\fbox{
\parbox{0.92\textwidth}{
\centering
parse\_query $\rightarrow$ retrieve\_data $\rightarrow$ extract\_evidence
$\rightarrow$ analysis\_agent $\rightarrow$ verification\_agent
$\leftrightarrow$ revise\_sections $\rightarrow$ prediction\_agent
$\rightarrow$ report\_writer
}}
\caption{High-level ERRGen pipeline.}
\label{fig:pipeline}
\end{figure}

The pipeline is implemented as a typed state graph. In the full mode, the
analysis sections go through verification and revision before the prediction
section is generated. In the baseline mode, the investment recommendation is
generated directly without the verification-revision loop, which makes it a
useful ablation for evaluation.

\subsection{Data Collection and Evidence Construction}
ERRGen retrieves real data from multiple sources, including Financial Modeling
Prep (FMP), NewsAPI, Finnhub, SEC filings, and Yahoo Finance. The use of real
APIs ensures that the system operates on actual external evidence rather than
synthetic or manually curated examples.

All retrieved content is normalized into structured evidence units together
with source metadata. This common evidence representation allows later stages
of the pipeline to reason over heterogeneous sources in a unified manner.

You can describe the role of each source here:
\begin{itemize}[leftmargin=1.5em]
    \item FMP for company profile, financial statements, stock news, and price
    history;
    \item NewsAPI for supplemental company news;
    \item Finnhub for optional fallback news and price history;
    \item SEC for company filings and financial fallback data;
    \item Yahoo Finance for price history fallback.
\end{itemize}

\subsection{Evidence Extraction and Deterministic Financial Calculation}
After retrieval, ERRGen performs two extraction workflows. The first is a
deterministic financial extraction process that identifies structured numeric
facts and creates calculation requests. The second is a news extraction process
that uses a language model to summarize events and signals from news articles.

A central design choice is that the language model never performs free-form
arithmetic for financial metrics. Instead, all numerical reasoning is delegated
to a deterministic calculator that supports operations such as growth rate,
margin, current ratio, debt-to-equity, net debt, free cash flow margin, and
company-versus-benchmark price return. This improves reproducibility and makes
each numeric claim auditable.

\subsection{Section-Level Report Generation}
The report is generated section by section instead of being produced in a
single pass. The current system includes the following major analytical
sections:
\begin{itemize}[leftmargin=1.5em]
    \item Company Overview
    \item Recent Developments
    \item Financial Analysis
    \item Business \& Competitive Analysis
    \item Risk Analysis
    \item Investment Recommendation \& Outlook
\end{itemize}

Different sections receive different subsets of evidence and calculations. For
example, recent developments rely more on news evidence, while financial
analysis relies more on statement data and deterministic metrics. This design
reduces irrelevant context and encourages more focused section-level reasoning.

In addition to narrative sections, the final report also includes deterministic
summary tables such as company snapshot, key financials, key ratios, recent
filings, and market performance when the underlying evidence is available.

\subsection{Paragraph-Level Citation Tracking}
Each generated paragraph stores explicit references to the evidence chunks and
calculation results that support it. This paragraph-level citation design makes
the final report traceable at a finer granularity than document-level source
lists. It also allows the verification system to validate each paragraph
against the exact evidence it cites.

\subsection{Verification and Revision Loop}
The verification mechanism is one of the core components of ERRGen. Every
generated paragraph begins as an unverified draft. A checker agent then
evaluates whether the paragraph is adequately supported by its cited evidence
and calculations. If problems are found, a reviser agent receives the checker
feedback and produces a corrected paragraph. The revised paragraph is then
checked again.

The checker currently evaluates seven issue types:
\begin{itemize}[leftmargin=1.5em]
    \item unsupported claim
    \item citation mismatch
    \item hallucination
    \item numerical error
    \item internal inconsistency
    \item scope violation
    \item overclaiming
\end{itemize}

The loop terminates when the paragraph passes verification or when the maximum
revision limit is reached, in which case the paragraph is marked as unresolved.
This design improves reliability while guaranteeing that the workflow finishes.

\subsection{Prediction Gate}
The investment recommendation section is generated only when upstream sections
do not contain unresolved major or critical issues. If such issues remain, the
prediction step is skipped and a warning is inserted into the report. This
prevents the system from producing a strong investment conclusion on top of
unreliable analysis.

\subsection{Run Artifacts and Auditability}
ERRGen stores a full set of run artifacts, including the parsed request,
retrieved sources, evidence chunks, extracted facts, calculation requests,
calculation results, paragraph drafts, checker verdicts, revision records, and
final report output. This makes the system easier to debug, evaluate, and
analyze as a research prototype.

\subsection{Experimental Setup}
Our experiments evaluate ERRGen along two main axes. First, we compare the
full system against an ablation baseline that generates reports without the
verification-revision loop. Second, we compare full ERRGen reports against
ground-truth human-written reports.

The evaluation pipeline uses an LLM-as-a-judge framework. It first applies an
accuracy-oriented evaluation, including claim support, numerical accuracy,
citation precision, temporal validity, and consistency. If a report passes the
accuracy gate, it is then scored on quality-oriented dimensions such as
financial analysis quality, news coverage, company-market-industry discussion,
investment usefulness, risk discussion quality, and writing quality.

\subsection{Quantitative Results}
Table~\ref{tab:main-results} summarizes the current evaluation results.

\begin{table}[H]
\centering
\begin{tabular}{lcccc}
\toprule
Experiment & Samples & Wins & Losses & Ties \\
\midrule
Full vs Baseline & 20 & 12 & 4 & 4 \\
Full vs Ground Truth & 10 & 0 & 10 & 0 \\
\bottomrule
\end{tabular}
\caption{Summary of current pairwise evaluation results.}
\label{tab:main-results}
\end{table}

For the full versus baseline comparison, the full ERRGen system achieved an
adjusted win rate of 0.70. This indicates that adding verification, revision,
and gating improves the overall quality and reliability of the generated
reports relative to an unverified generation pipeline.

However, the comparison against ground-truth reports shows that the generated
outputs still trail human-written equity research in depth, synthesis, and
overall investor usefulness. This suggests that factual control alone is not
sufficient to match expert-level report quality.

\subsection{Qualitative Analysis}
This subsection should include one or two representative cases.

For a successful case, discuss a report that passed verification and produced a
complete recommendation section. Explain how evidence grounding, deterministic
metrics, and the verification loop contributed to a more reliable final output.

For a failure case, discuss a report in which one or more sections remained
unresolved and the recommendation was skipped. Explain the underlying causes,
such as insufficient qualitative evidence, noisy news retrieval, incomplete
filing support, or revision failure.

\subsection{Discussion}
The experimental results suggest that the verification-oriented design improves
report reliability in comparison with an unverified baseline. At the same
time, the gap relative to human-written reports highlights that financial
analysis quality depends not only on factual correctness but also on deeper
synthesis, stronger valuation reasoning, and better use of domain knowledge.

\section{Conclusion, Discussion and Future Work}
This project presented ERRGen, an evidence-grounded equity research report
generator designed to improve factual reliability, numerical correctness, and
traceability in long-form financial text generation. The system combines real
data retrieval, structured evidence representation, deterministic financial
calculation, section-level generation, paragraph-level citation tracking, and a
checker-reviser verification loop.

The results show that the full system outperforms an unverified baseline,
demonstrating the practical value of evidence grounding and iterative
verification. However, the project also reveals that matching the depth and
decision usefulness of expert-written research reports remains challenging.

Several directions are promising for future work. First, the system could be
extended with a valuation module, including discounted cash flow and multiple
based analysis. Second, richer filing parsing and better industry-level data
could improve business and risk discussion quality. Third, stronger retrieval
and reranking methods may improve evidence coverage. Finally, combining
automatic evaluation with human expert assessment would provide a more complete
understanding of the system's strengths and weaknesses.

\newpage
\bibliographystyle{plain}
\bibliography{references}

\appendix
\section{Appendix (Optional)}
You can place supplementary material here, such as:
\begin{itemize}[leftmargin=1.5em]
    \item prompt templates;
    \item schema definitions;
    \item additional experiment tables;
    \item extra case studies;
    \item screenshots or example run artifacts.
\end{itemize}

\end{document}
